\chapter{Introduction}
\label{chap:introduction}

\section{Words draw boundary lines across the ocean of ``concepts''}
When I was still in elementary school, my favourite biography in the library was Helen Keller's. Stricken by illness as a child, she lost both sight and hearing and bore the ``triple hardship'' of being blind, deaf, and mute, yet devoted herself to the advancement of education and welfare for the disadvantaged.

But Helen Keller's story is not merely a tale of ``overcoming adversity.'' What is etched most vividly in my memory is the moment when she first understood the word ``water.'' At first, Helen insisted that ``cup'' and ``the water in it'' were the same thing. However, the instant when Anne Sullivan signed ``W-A-T-E-R'' into her hand while letting the cold water from a pump run across Helen's palm, her world changed. When she could separate the object that flows from the container, the continuum of chaos was reshaped into an ordered world with partitions.

This scene exemplifies the essence of language. We do not grasp the external world directly as ``things-in-themselves.'' The information flowing into our senses is, in truth, a continuum; language draws borders there and carves out units of meaning. Even the boundary between river and sea is not sharply inscribed by nature itself; it is defined for human convenience and by culture~-- ``this far is river,'' ``from here is sea.'' Words do not merely copy lines already inscribed upon the world; they are the very act of drawing those lines as needed. In other words, language is not just a label: it is the activity by which we impose boundaries within our cognition and shape the world we see. To learn words is not merely to expand vocabulary; it is to gain a way of partitioning the ocean of concepts and choosing how to look at the world. In that sense, we are, day by day, sketching new boundaries.

\section{Verbalisation is a lossy compression}
Verbalisation is an act of forcing the complex and continuous sensation into discretised vocabulary and grammar. When we perform \textbf{verbalisation}, there lurks a commonly overlooked risk; it is intrinsically \textbf{lossy}. The moment we verbalise, the fine nuances contained within raw feeling are abstracted into the fences of cognition that humanity has drawn at the level of words. Even if we choose the same term, the warmth that term carries depends on each person's embodied history, and it is eventually averaged out. The relief or pride we felt when we opened our old school yearbook for the first time in a decade, once put into text, may be flattened into a bland ``how nostalgic.'' Moreover, lexicons and meanings for emotions differ greatly across languages and cultures\cite{jackson2019emotion}; terms like ``anger'' or ``joy'' do not necessarily point to the same nuances around the world. And yet humanity cannot let go of language. Even if it is a lossy compression, language remains the only standardised interface that connects the human heart to society.


\section{Objectives}
The goal of this project is not to deny language but to design a parallel circuit that supplements the components lost in the compression of verbalisation, i.e.\ a circuit for restoring what words leave behind. What if we could scoop up the human central cognition directly, without passing through the rendering pipeline of language, and translate it into other forms of expression? Could we transfer the lingering sensation that arises when reading a letter from your significant other into poetry or essays? Or could we improve novel translations by capturing emotions lost through differences in lexical definitions between languages? Here lies an uncharted gap in our understanding of human cognition. 

This project aims to emulate a kind of human synaesthesia. In other words, we aim to map emotions across modalities. We will develop the foundational technology to extract, manipulate, and re-synthesise a \textbf{pre-verbal, intermediate representation of emotion}, and explore possible ways to steer the model output using those representations of emotion.

The final product of this project is a computational system capable of extracting, manipulating, and presenting the emotions contained in input text without having to go through the compression of verbalisation.
The system focuses not on what the text superficially tells, but on the emotional state that existed immediately before the text was produced, and treats it as a continuous vector representation.

When a user enters any text, the system analyses the emotional nuances contained in that text and visualises them as a low-dimensional continuous vector. These emotion vectors are not discrete labels such as ``joy'' and ``sadness,'' but are expressed as multiple overlapping emotional directions, thus preserving subtle mixed emotions and fluctuations that cannot be captured by language.

Users can intuitively manipulate this emotion code space. For example, by strengthening only a specific emotional direction and weakening another direction for the simultaneous presence of \texttt{relief} and \texttt{loneliness} in a certain sentence, it is possible to sensitively edit the emotional state that one really wanted to convey. It is important to note that this operation is not performed at the linguistic level, such as vocabulary selection or stylistic adjustment, but rather on the emotional expression itself before the language is generated.

The edited emotion code is again transformed into a linguistic expression. The text generated here is not simply a paraphrased sentence, but is the result of a re-expression that retains or intentionally transforms the emotion contained in the original sentence.

In this process, the user does not need to put their emotions into explicit words. The verbal averaging and simplification that is usually considered inevitable between emotion recognition and generation is bypassed, and the raw emotion felt by the user can be transferred to another form of expression in a more controllable manner.

This end product is neither a system for classifying and diagnosing emotions nor a tool for automatically generating emotional sentences. Rather, it is an interface that temporarily holds the emotions that people had before they put them into words as manipulatable intermediate expressions, from which new linguistic expressions can be launched. The ultimate value of this project lies in the fact that it makes it possible to return to the boundaries drawn by language and choose words again.


%%%%%%%%%%%%%%%%%%%%%%%%%%%%%%%%%%%%